\documentclass[a4paper,11pt]{article}
\usepackage[a4paper,top=1.6in,bottom=1.6in,left=1.3in,right=1.3in]{geometry}
\usepackage[full]{textcomp}
\usepackage[osf]{newtxtext}

\usepackage{microtype}

\usepackage{amssymb}             % for blackboard bold, etc
\usepackage{mathtools}

\usepackage{tikz}
\usepackage{amsthm}                % better theorem environment
\usepackage{mhequ}
\usepackage{hyperref}

\newcommand{\eqdef}{\stackrel{\mbox{\tiny def}}{=}}

\numberwithin{equation}{section}

% various theorems, numbered by section
\newtheorem{lemma}{Lemma}[section]
\newtheorem{proposition}[lemma]{Proposition}


\DeclareMathOperator{\id}{id}
\def\Wick#1{\mathopen{:}#1\mathclose{:}}


\newcommand{\EE}{\mathbb{E}}     %mathematical expectation
\newcommand{\RR}{\mathbb{R}}      % for Real numbers
\newcommand{\TT}{\mathbb{T}}

\newcommand{\cC}{\mathcal{C}}
\newcommand{\dD}{\mathcal{D}}
%\def\CO{\mathcal{O}}

%\newcommand{\Poly}{\text{\tiny Poly}}


\let\f\frac

\def\d{\partial}


\newcommand{\eps}{\varepsilon}

\colorlet{symbols}{blue!90!black}
\colorlet{testcolor}{green!60!black}

\tikzset{
	smalldot/.style={circle,fill=symbols,draw=symbols, solid,inner sep=0pt,minimum size=0.5mm},
}

\def\scal#1{\langle#1\rangle}
\def\cent#1{\mathopen{{\langle\kern-0.3em\rangle}}#1\mathclose{{\langle\kern-0.3em\rangle}}}
\def\bscal#1{\big\langle#1\big\rangle}
\def\Bscal#1{\Big\langle#1\Big\rangle}


\makeatletter
\def\DeclareSymbol#1#2#3{\expandafter\gdef\csname MH@symb@#1\endcsname{\tikz[baseline=#2,scale=0.15,draw=symbols]{#3}}\expandafter\gdef\csname MH@symb@#1s\endcsname{\scalebox{0.65}{\tikz[baseline=#2,scale=0.15,draw=symbols]{#3}}}}
\def\<#1>{\csname MH@symb@#1\endcsname}
\makeatother

\DeclareSymbol{0}{0}{\draw[white] (-.4,0) -- (.4,0); \draw (0,0)  -- (0,1.8) node[smalldot] {};}
\DeclareSymbol{1}{-1}{\draw[white] (-.4,0) -- (.4,0); \draw (0,-0.5)  -- (0,1.3) node[smalldot] {};}
\DeclareSymbol{2}{-1}{\draw (-0.7,1.3) node[smalldot] {} -- (0,-0.5) -- (0.7,1.3) node[smalldot] {};}
\DeclareSymbol{3}{1}{\draw (-1,1.5) node[smalldot] {} -- (0,0) -- (1,1.5) node[smalldot] {}; \draw (0,2) node[smalldot] {} -- (0,0);}
\DeclareSymbol{20}{-.2}{\draw (-0.7,1.5) node[smalldot] {} -- (0,0.5); \draw (0.7,1.5) node[smalldot] {} -- (0,0.5) -- (0,-0.5);}
\DeclareSymbol{30}{0}{\draw (-1,1.5) node[smalldot] {} -- (0,0.5) -- (1,1.5) node[smalldot] {}; \draw (0,2) node[smalldot] {} -- (0,0.5) -- (0,-.5);}
\DeclareSymbol{31}{0}{\draw (-1,1.5) node[smalldot] {} -- (0,0.5) -- (1,1.5) node[smalldot] {}; \draw (0,2) node[smalldot] {} -- (0,0.5) -- (0,-.5) -- (1,0.5) node[smalldot] {} ;}


\begin{document}

\title{(Lack of) quasi-shift invariance of the $\Phi^4_3$ measure}
\author{Martin~Hairer}
\date{}
\maketitle

\section{Question}


Let $\TT^3$ be the three dimensional unit size torus and let $\mu$ be the $\Phi^4_3$ measure on the space of distributions $\dD'(\TT^3)$. Let $\psi : \TT^3 \to \RR$ be a smooth function that is not identically zero and let $T_\psi : \dD'(\TT^3) \to \dD'(\TT^3)$ be the shift map given by $T_\psi(u) = u + \psi$ (with the usual identification of smooth functions as distributions). Is the statement ``the measures $\mu$ and $T_\psi^* \mu$ are equivalent'' true? Here, equivalence of measures is in the sense of having the same null sets and $T_\psi^*$ denotes the pushforward under $T_\psi$.


\section{Some Context}

One of the very few interacting quantum field theories that can be rigorously constructed is the
so-called (bosonic) $\Phi^4$ theory in (space-time) dimensions $2$ and $3$. It has long been known that
in dimension $2$ and finite volume there is a natural identification between the Hilbert space of the
interacting theory and that of the corresponding free theory. On the other hand, 
Glimm \cite{Glimm} observed that this is no longer the case in dimension $3$. At the level of the 
corresponding Euclidean theories (which are represented by probability measures on the
space of Schwartz distributions on the corresponding space-time), this translates into the fact that the 
$\Phi^4$ measure $\mu$ and the corresponding free field measure $\nu$ are equivalent in dimension $2$
but mutually singular in dimension $3$. In fact, there is a sense in which the dimension that delimits
between the two behaviors is $8/3$. It is then natural to ask in which dimensions $\mu$
has the weaker property that $\mu$ and $T_\psi^* \mu$ are equivalent for smooth $\psi$. Here it turns out
that the borderline dimension is $3$, and the question probes on which side it falls.

\subsection{An incomplete heuristic}

Regarding the proof, a tempting heuristic is to use the fact that one should think of $\mu$ as having the
density with respect to ``Lebesgue measure on $\dD'$'' (which of course doesn't exist) proportional to
\begin{equ}
\exp \Bigl(- \int_{\TT^3} \Bigl(\frac12 |\nabla\Phi(x)|^2 + \frac 14|\Phi(x)|^4 - \frac C2 |\Phi(x)|^2\Bigr)\,dx\Bigr)\;,
\end{equ}
where $C$ is a (diverging) constant of the form $C = 3c_1 - 9 c_2$, where
$c_1$ is the expectation of $|\Phi(x)|^2$ under the free field measure $\nu$ (which is of course infinite)
and $c_2$ is an additional logarithmically divergent constant. The density of $T_\psi^* \mu$
with respect to $\mu$ is then formally given by
\begin{equs}
\exp \Bigl(- \int_{\TT^3} &\Bigl(\frac12 |\nabla\psi(x)|^2 + \frac 14|\psi(x)|^4 + \Phi(x) \Delta \psi(x)
- \Phi(x)\psi^3(x) \\
& - \psi(x) \bigl(\Phi^3(x) - C\Phi(x)\bigr) + \frac{\psi^2(x)}{2} \bigl(3\Phi^2(x) - C\bigr)\Bigr)\,dx\Bigr)\;,
\end{equs}
Since the terms on the first line are well-defined for smooth $\psi$ and 
one expects $\Phi^3 - C\Phi$ and $\Phi^2 - c_1$ to be quite well-behaved, the additional
logarithmically divergent term proportional to $c_2$ causes this ``density'' to diverge, suggesting
(correctly) that $\mu$ and $T_\psi^* \mu$ are mutually singular.

There are at least two problems with such an approach. First, $\Phi^3 - C\Phi$ 
does actually \textit{not} define a random distribution, whether $\Phi$ is distributed according
to $\mu$ or to the free field $\nu$ (which guides the heuristic). This is because if it were,
it would have a covariance behaving like $|x-y|^{-3}$ around the diagonal, which is not integrable 
in dimension $3$. The second problem is that such an argument suggests that, if 
$\mu_n = \exp(- f_n)\,\nu$ for some ``nice'' probability measure $\nu$ and functions $f_n$ that 
fail to converge to a ``nice'' limit, then $\mu_n$ fails to converge to a limit $\mu$. This 
of course is not true: for a suitable (diverging) sequence of constants $c_n$, the sequence
$f_n(x) = c_n + n \cos(nx)$ is such that if $\nu$ is Lebesgue measure on $[0,1]$, then $\mu_n$ 
converges weakly to Lebesgue measure even though the log-densities $f_n$ fail to converge.
Any proof needs to be based on a different approach or to satisfactorily address these problems.

\subsection*{Acknowledgement}

{\small
The proof presented below is a simplified version of an argument that will appear as part of 
a joint publication with Jacopo~Peroni. This in turn is based on the article \cite{HKN24}.
}

\section{Notations}

We fix a space-time white noise $\xi$ on $\RR \times \TT^3$.
We define \<1> as the stationary solution to the linear equation
    \[
        (\d_t + 1-\Delta)\<1> = \xi, \qquad \mbox{on } \RR \times \TT^3\;.
    \]
    (We use the convention that symbols represent random space-time distributions rather than elements of a regularity structure.)
Starting from this process, we define \<2> and \<3> as its Wick square and cube respectively, which are given by
\begin{equ}
\<2> = \lim_{N \to \infty} H_2(\<1>_N,c_{N})\;,\qquad
\<3> = \lim_{N \to \infty} H_3(\<1>_N,c_{N})\;,
\end{equ}
where $\<1>_N = P_N \<1>$ and $c_{N} = \EE \<1>_N^2$ (which is constant in space and time).
Here, $P_N$ denotes the projection onto Fourier modes with $|k| \le N$ and $H_n$ denotes the $n$th
Hermite polynomial normalised such that $H_0 \equiv 1$, $H_n' = nH_{n-1}$, and 
$\EE H_n(Z,1) = 0$ for a normal random variable $Z$. 
The first convergence takes place in the space of continuous functions of time with values
in $\cC^{-1-2\kappa}$, while the second convergence takes place in the space-time parabolic 
space $\cC^{-\frac32-3\kappa}$.

With these notations in place, we define \<20> as the stationary solution to 
    \[
        (\d_t + 1-\Delta)\<20> = \<2>\;,
    \]
    and similarly for \<30>.
    For a more comprehensive and pedagogical introduction to the general tree-like notation, we refer the reader to \cite{MWX15}.
    We also write ``$\prec$'' for Bony's paraproduct (in space) as defined for example in \cite[Sec.~2.1]{GIP12}
    and, given a random $N$-dependent process $w$, we will sometimes use the physicists shorthand notation $\Wick{w^k}$
    instead of $H_k(w,c_{N})$.


\section{Answer and Proof}


The statement is \textbf{false}. In particular, for any smooth function $\psi \not \equiv 0$ and any choice of the
parameters involved in the definition of $\mu$ (mass and coupling constant, provided that the latter is non-zero),
the measures $\mu$ and $T_\psi^* \mu$ are  mutually singular. 

For notational simplicity, we fix the mass and the
coupling constant to $1$, but this has no incidence on the proof.
Our main starting point is the following statement, a proof of which 
can be found for example in \cite{Konstantin} and \cite[Lemma 4.19]{EW24} for \eqref{e:processu}, combined with \cite{CC18}
for \eqref{e:paracontrolled} (see Ansatz~2.11 there). 
Throughout this proof, $\kappa>0$ is chosen small enough 
($\kappa = 1/100$ is certainly sufficient).

\begin{proposition} \label{v-regularity}
    There exists a stationary process $v$ that is almost surely continuous
    in time with values in $\cC^{1-2\kappa}(\TT^d)$ and such that the process
    \begin{equ}[e:processu]
        u = \<1> - \<30> + v\;,
    \end{equ}
	is stationary with fixed time distribution equal to $\mu$. Furthermore, the process $v$ is such that
    \begin{equ}[e:paracontrolled]
        v = -3 \bigl((v -\<30>) \prec \<20>\bigr) + v^\sharp\;,
    \end{equ}
    where $v^\sharp$ is continuous with values in $\cC^{1+4\kappa}(\TT^d)$. 
\end{proposition}


It was furthermore shown in \cite[Lemmas~3.1 \& 3.4]{HKN24} (but see \cite{Hai14a} for a similar
result using a slightly different regularisation) that the 
processes $\<1>$ and $\<30>$ are almost surely continuous in time with values in 
$\cC^{-\frac12-\kappa}$ and $\cC^{\frac12-3\kappa}$ respectively.

Before we proceed, we remind some notations and preliminary results.
First of all, we define the additional diverging constant
\begin{equation}
    c_{N,2} := \EE \left [ \<2>_N \<20>_N \right ],
\end{equation}
where $\<2>_N:=P_N\<2>$ and $\<20>_N := P_N\<20>$.
The main ingredient of our proof is the event
\begin{equ}
    B^{\gamma}:= \left \{ u \in \dD' : \lim_{N \to \infty}\left \langle (\log N)^{-\gamma} \left ( H_3 \left ( P_N u; c_{N} \right ) + 9c_{N,2} P_Nu\right ),\psi \right \rangle_{\TT^3} = 0 \right \},
\end{equ}
which will be used to distinguish between the shifted and the non-shifted measures.

We will also use the following two technical lemmas whose proofs can be found in Section~\ref{proof-lemma-section}. These are very similar to \cite[Lemma 3.11 and Lemma 3.12]{HKN24}.
\begin{lemma} \label{Z_to_zero_sigma}
Let  $\gamma > \frac12$. Then, for any fixed $t>0$,
        \begin{equ}
            \lim_{N\to\infty} (\log N) ^{-\gamma} \Wick{\left (  \<1>_N \right )^3}(t)  = 0
        \end{equ}
        almost surely in $\cC^{-\frac{3}{2}}\bigl ( \TT^{3} \bigr )$ and in $L^p \bigl ( \Omega; \cC^{-\frac{3}{2}} \bigl ( \TT^3 \bigr ) \bigr )$ for any $p >0$.
\end{lemma}
\begin{lemma} \label{cN2_div_frac}
    For $N$ large, one has   $ c_{N,2} \gtrsim 
            \log N $.
\end{lemma}

The following results are essentially standard, but we recall their statements for later reference.

\begin{lemma}\label{lem:convSimple}
For any polynomial $P$, the expression $\<1>_N P(\<30>_N)$ converges almost surely to some finite limit in 
$\cC^{-\frac12 -\kappa}$.
\end{lemma}

\begin{proof}
By paralinearisation and standard commutator estimates (see \cite[Lems~2.4 \& 2.6]{GIP12}) it suffices to 
consider the case $P(x) = x$. This is by now standard, see for example \cite[Sec.~4.4]{CC18}.
\end{proof}

\begin{lemma}\label{lem:convParacontrolled}
Let $v$ be a process satisfying the decomposition \eqref{e:paracontrolled}. Then, the expressions
$\Wick{\<1>^2_N}\<30>_N - 3c_{N,2} \<1>_N$ and $\Wick{\<1>_N^2}v_N + 3c_{N,2} (v_N - \<30>_N)$ both converge
almost surely to finite limits in $\cC^{-1 -2\kappa}$ as $N \to \infty$.
\end{lemma}

\begin{proof}
Regarding the first expression, its convergence was essentially for example in \cite[Sec.~4.6]{CC18}. (The
approximation used there is slightly different, but the differences are unimportant.)
Regarding the second expression, the claim follows from \cite[Sec.~4.5]{CC18} (modulo again unimportant changes in 
the approximation scheme), combined with the commutator estimate \cite[Lem.~2.4]{GIP12}.
\end{proof}

We now turn to the proof of the main claim. For this, we first claim that if $u$ is as in \eqref{e:processu},
then, for any fixed $t$, one has $u(t) \in B^\gamma$.
Indeed, writing $u_N$ as a shorthand for $P_N u$ and expanding the Wick power, we have
    \begin{align*}
        (\log N)^{-\gamma} &H_3 \left ( u_N; c_{N} \right )
        =(\log N)^{-\gamma} H_3 \left (  \left ( \<1>_N - \<30>_N + v_N \right ); c_{N} \right )\\
        &=(\log N)^{-\gamma} \sum_{i=0}^3 \binom{3}{i} \Wick{\left (  \<1>_N \right )^i} \left (- \<30>_N + v_N \right )^{3-i}\\
        &=(\log N)^{-\gamma} \sum_{i=0}^3 \sum_{j=0}^{3-i} \binom{3}{i}\binom{3-i}{j} \Wick{\left (  \<1>_N \right )^i} \left ( - \<30>_N \right )^{j}\left ( v_N \right )^{3-i-j}\\
        &=(\log N)^{-\gamma}\Wick{ \<1>_N^3}
        - 3(\log N)^{-\gamma}\Wick{\<1>_N^2} \, \<30>_N  + 3(\log N)^{-\gamma}\Wick{\<1>_N^2}  \, v_N \\
        &\quad+ (\log N)^{-\gamma} \sum_{\substack{0 \le i+j \le 3 \\ (i,j) \neq (3,0),(2,1),(2,0)}} \binom{3}{i}\binom{3-i}{j} \Wick{\left (  \<1>_N \right )^i} \left ( - \<30>_N \right )^{j}\left ( v_N \right )^{3-i-j}.
    \end{align*}
    The first term $(\log N)^{-\gamma}\Wick{\<1>_N^3}$ and the terms present in the last sum all converge to $0$ by Lemma~\ref{Z_to_zero_sigma} (given that $\gamma > \frac{1}{2}$), standard product estimates 
(e.g. \cite[Theorem 2.5]{Bon81} or \cite[Proposition 2.3]{MWX15}) and Lemma~\ref{lem:convSimple}.
    
    It therefore remains to show that $-\Wick{\<1>^2_N}\<30>_N + \Wick{\<1>_N^2}v_N + 3c_{N,2} u_N$ also
converges to zero almost surely in the sense of distributions. We rewrite this term as
\begin{equ}
\Wick{\<1>_N^2}v_N + 3c_{N,2} (v_N - \<30>_N) -\bigl(\Wick{\<1>^2_N}\<30>_N - 3c_{N,2} \<1>_N\bigr)\;.
\end{equ}
By Lemma~\ref{lem:convParacontrolled} we know that this expression converges to an element of $\cC^{-1-2\kappa}(\TT^d)$, whence we conclude that
    \[
        \bigl\langle\left ( \log N \right )^{-\gamma} \left ( -\Wick{\<1>^2_N}\<30>_N + \Wick{\<1>_N^2}v_N + 3c_{N,2} u_N \right ),\psi\bigr\rangle \xrightarrow{N \to \infty} 0
    \]
    almost surely, thus proving that $\mu(B^{\gamma}) = 1$.
    
    In order to conclude the proof, it suffices to show that $u + \psi \not \in B^\gamma$. 
For this, we expand similarly to before the expression appearing in this event as
    \begin{align*}
        (\log N)^{-\gamma}&H_3 \left( \left ( u_N + \psi_N \right ); c_{N} \right ) + (\log N)^{-\gamma}9c_{N,2} \left ( u_N + \psi_N \right )\\
        &=(\log N)^{-\gamma} \sum_{i=0}^3 \binom{3}{i} \Wick{\left ( u_N \right )^i} \left ( \psi_N \right )^{3-i} + (\log N)^{-\gamma}9c_{N,2} \left ( u_N + \psi_N \right )\\
        &=(\log N)^{-\gamma}\Wick{\left ( u_N \right )^3} + (\log N)^{-\gamma}9c_{N,2} u_N\\
         &\qquad + 3(\log N)^{-\gamma} \Wick{\left ( u_N \right )^2} \left ( \psi_N \right ) + 3(\log N)^{-\gamma} \left ( u_N \right ) \left ( \psi_N \right )^2\\
         &\qquad+ (\log N)^{-\gamma}\left ( \psi_N \right )^3 + (\log N)^{-\gamma}9c_{N,2} \psi_N.
    \end{align*}
    The sum of the first two terms was just shown to converge to $0$ almost surely in $\cC^{-\frac{3}{2}-3\kappa}(\TT^d)$ for $N\to \infty$.
    
Since $\Wick{u_N^2}$ and $u_N$ both converge to finite distributional limits almost surely by Lemma~\ref{lem:convSimple}, 
the next three terms also converge to $0$ almost surely.
    
    Concerning the last element however, we know from Lemma \ref{cN2_div_frac} that
    \[
        \left ( \log N\right )^{-\gamma} c_{N,2} \gtrsim \left ( \log N\right )^{-\gamma+1}\;.
    \]
Since the contribution of this term to the expression in the event $B^\gamma$ is given by
\begin{equ}
9\left ( \log N\right )^{-\gamma} c_{N,2} \scal{\psi_N,\psi}\;,
\end{equ}
and since $\scal{\psi_N,\psi} \to \|\psi\|^2 > 0$, this diverges, whence we conclude that 
$u + \psi \not \in B^\gamma$ and therefore $\bigl(T_\psi^* \mu\bigr)(B^\gamma) = 0$, so that 
$T_\psi^* \mu$ and $\mu$ are mutually singular.



\subsection{Proof of the lemmas} \label{proof-lemma-section}

\begin{proof}[Proof of Lemma \ref{Z_to_zero_sigma}]
    We use the embedding $W^{\beta, 2p} \hookrightarrow W^{\beta - \frac{d}{2p}, \infty} \hookrightarrow \cC^{\beta - \frac{d}{2p}}$, with $\beta = -\frac{d}{2}$.
    Using the definition of $W^{\beta,2p}$ norm and the equivalence of moments for Gaussian polynomials, one has
    \begin{align*}
        \EE \left [ \left \lVert (\log N)^{-\gamma} \Wick{\left ( \<1>_N \right )^J} \right \rVert^{2p}_{W^{-\frac{3}{2}, 2p}} \right ] \lesssim \int_{\TT^d} \EE \left [ (\log N)^{-2\gamma} \left | \left \langle \nabla \right \rangle^{-\frac{3}{2}} \Wick{\left ( \<1>_N \right )^3} \right |^{2} \right ]^p dx\;.
    \end{align*}
 Since one has
    \begin{align*}
        \EE \left [ (\log N)^{-2\gamma} \left | \left \langle \nabla \right \rangle^{-\frac{3}{2}} \Wick{\left ( \<1>_N \right )^3} \right |^{2} \right ]
        &\lesssim (\log N)^{-2\gamma} \sum_{|\omega_i| \le N} \left \langle \omega_1 + \dots + \omega_3 \right \rangle^{-3} \prod_{i=1}^3 \left \langle \omega_i \right \rangle^{-2}\\
        &\lesssim (\log N)^{-2\gamma} \sum_{r_1=0}^N \frac{r_1^{2}}{\left ( 1+r_1^2\right )^{\frac{5}{2}}} r_1^{2}
        \lesssim (\log N)^{-2\gamma + 1}\;,
    \end{align*}
the desired result follows from a standard Borel--Cantelli argument.
\end{proof}
Next, we prove Lemma~\ref{cN2_div_frac}, which provides a lower bound on the parameter $\gamma$. This bound ensures that the event $A^\gamma$ (or $B^\gamma$) is distinguishable under the shifted measure when compared to the non-shifted one.
\begin{proof}[Proof of Lemma \ref{cN2_div_frac}]
    Expanding the definition of $c_{N,2} := \EE \left [ \<2>_N \<20>_N \right ]$, we get
    \begin{equs}
	    c_{N,2} &= 2\sum_{\substack{\omega_1 + \omega_2 = \omega_3 \\ |\omega_i| \le N}} \int_{\RR} \hat{P}_{t-u}(\omega_3) \int_{\RR} \hat{P}_{t-u_1}(\omega_1)\hat{P}_{u-u_1}(-\omega_1)du_1\\
        \quad&\times\int_{\RR}\hat{P}_{t-u_2}(\omega_2)\hat{P}_{u-u_2}(-\omega_2)du_2du\\
        &\simeq \sum_{\substack{\omega_1 + \omega_2 = \omega_3 \\ |\omega_i| \le N}} \int_{\RR} e^{-|t-u|\langle \omega_3 \rangle^2} \frac{e^{-|t-u|\langle \omega_1 \rangle^2}}{\langle \omega_1 \rangle^{2}} \frac{e^{-|t-u|\langle \omega_2 \rangle^2}}{\langle \omega_2 \rangle^{2}}du\\
        & \gtrsim \sum_{\substack{|\omega_i| \le N}} \frac{1}{\langle \omega_1 \rangle^{2}}\frac{1}{\langle \omega_2 \rangle^{2}}\frac{1}{\langle \omega_1 \rangle^2+\langle \omega_2 \rangle^2+\langle \omega_1+\omega_2 \rangle^2}\\
        & \gtrsim \sum_{\substack{|\omega_i| \le N}} \frac{1}{1 + |\omega_1|^{2}}\frac{1}{1 + |\omega_2|^{2}}\frac{1}{1+|\omega_1|^{2} \vee |\omega_2|^{2}}\\
        & \gtrsim \sum_{\substack{|\omega_1| \le |\omega_2| \le N}} \frac{1}{1 + |\omega_1|^{2}}\frac{1}{1 + |\omega_2|^{4}}\;.
    \end{equs}
Bounding the sum by an integral, we finally conclude that this expression is bounded
from below by a multiple of
    \begin{equ}
       \int_0^N \frac{r^2}{1+r^2} \int_r^\infty \frac{s^2}{1+s^4}\,ds\,dr
       \gtrsim \int_0^N  \frac{1}{1+r^2}\,dr \simeq \log N\;,
      \end{equ}
     as claimed.
\end{proof}

\bibliographystyle{alpha}
\bibliography{Refs}


\end{document}

\section{Comments}

We observed three different failure modes: 
\begin{enumerate}
\item The 
LLM believes that the $\Phi^4_3$ measure $\mu$ has a Radon--Nikodym derivative proportional
to $\exp \bigl(- \int \Wick{\Phi^4}\,dx\bigr)$ with respect to the Gaussian free field (GFF) measure.
It then (correctly) deduces that, by the Cameron--Martin theorem, the measures $\mu$ and
$T_\psi^*\mu$ are equivalent. The problem of course is that $\mu$ and the GFF are mutually singular
in dimension $3$.

\item The LLM knows that $\mu$ and the GFF are mutually singular
and that there is an additional mass renormalisation appearing in its approximations (the constant 
$c_{N,2}$ in the notes above). It then compares the formal densities of the original measure and
the shifted measure with respect to the GFF. This contains a linear combination of the terms 
$\int \Wick{\Phi^k} \psi^{4-k}$ with $k \in \{0,\ldots,3\}$, as well as a term proportional
to $c_{N,2} \int (2\Phi + \psi)\psi$. It then argues that the former are well-defined in the 
limit, while the latter diverges, so one expects mutual singularity. This is a good heuristic,
but it ignores the fact that $\Phi$ is distributed according to $\mu$, which is singular
with respect to the GFF, so it isn't clear a priori whether its Wick powers make sense
(in fact the third Wick power doesn't, even for the free field) and / or whether this could 
create additional diverging terms creating a cancellation. It also ignores the fact that it is 
perfectly possible to have a sequence of densities $f_N$ with respect to some reference measure $\nu$ that don't
converge to a limit pointwise (or that almost surely converge to $0$), but such that 
$f_N \nu$ nevertheless converges weakly to a limiting probability measures that is equivalent to $\nu$.
(Take $f_N(t) = c_N \exp(N \sin(Nt))$ with suitably chosen $c_N$ and Lebesgue on $[0,1]$ as
reference measure.)

\item The LLM finds the short note at \href{http://hairer.org/Phi4.pdf}{hairer.org/Phi4.pdf}
which states the correct result and gives an extremely short sketch of proof (it  
defines the set $B^\gamma$ and gives a hint of a few sentences on why it does the trick).
It then takes this at face value without attempting to turn it into a proper proof.
\end{enumerate}


